\section{Teoria}\label{Sec:Teoria}
\paragraph{Esempio di equazione.} %Paragrafo
Ok, quindi avete visto come si fa per generare una tabella o inserire una figura. E per quanto riguarda le Equazioni?\\
Semplice, nel modo seguente
\begin{equation}\label{eq:EinsteinField_eq}
R_{\mu \nu} - \dfrac{1}{2} R g_{\mu \nu} + \Lambda g_{\mu \mu} = \dfrac{8 \pi G}{c^4} T_{\mu \nu}
\end{equation}
E la si richiama nel testo con il solito comando {\it ref}, quindi l'Equazione di Campo di cui sopra è l'Eq.\,\ref{eq:EinsteinField_eq}.\\
Se ho necessità di scrivere più formule consecutivamente, si usa eqnarray:
\begin{eqnarray}
\dfrac{\partial \mathcal D}{\partial t} &=& \nabla \times \mathcal H \\
\dfrac{\partial \mathcal B}{\partial t} &=& -\nabla \times \mathcal E \\
\nabla \cdot \mathcal B  &=& 0 \\
\nabla \cdot \mathcal D  &=& 0 
\end{eqnarray}
E se invece di quattro le voglio considerare come fosse una sola:
\begin{eqnarray}\label{eq:Maxwell_eq}
\dfrac{\partial \mathcal D}{\partial t} &=& \nabla \times \mathcal H \nonumber \\
\dfrac{\partial \mathcal B}{\partial t} &=& -\nabla \times \mathcal E \nonumber \\
\nabla \cdot \mathcal B  &=& 0 \nonumber \\
\nabla \cdot \mathcal D  &=& 0 
\end{eqnarray}

% Come vedete, per questa sezione non ho creato un nuovo file Sezione_3.tex, ma l'ho direttamente definita qua sotto. A seconda di come più vi aggrada, potete fare un file per ogni sezione e richiamarlo nel file principale con \input{}, oppure mettere tutto nello stesso file. It's up to you.
\section{Conclusioni}\label{Sec:Conclusioni}
Buon lavoro.
